\section{Discussion}
\subsection{System Trade-offs}
The two systems offer different operating points. CG-ICS provides higher average accuracy and stronger small-target performance on the evaluated datasets, while INSID3 provides a much cheaper inference path. Reference sensitivity adds another dimension: CG-ICS improves worst-reference accuracy on both datasets but has greater reference variation on LVIS. Choosing a system therefore depends on the importance of accuracy, latency, resource availability, and consistency across demonstrations.

The encoder matrix and within-system ablations explain why the complete-system ranking cannot be reduced to a backbone ranking. DINOv3 leads the tested feature-processing matrix, yet CG-ICS achieves higher complete-system accuracy through a different combination of language, prompting, and decoding. INSID3 processing also improves every tested trunk. The representation and the procedure that converts it into a mask should therefore be evaluated together.

\subsection{Marginal Value of Search}
The strongest motivation for selective search is the difference between an average gain and an input-specific gain. With visual boxes, full search improves reference-averaged IoU for 194 of 400 COCO queries and 535 of 1,000 LVIS queries; it decreases IoU for 196 and 403, respectively. A positive dataset mean therefore does not imply that every query benefits. The observed search--prompt interaction further shows that the value of one component changes when another is already present.

A same-GPU sensitivity analysis retains H200 subsets of 100 queries per dataset. The interaction is $+0.66$ points on COCO (95\% CI $[-1.67,3.01]$) and $-3.47$ on LVIS ($[-6.57,-0.60]$). The LVIS direction persists, while the smaller COCO subset is inconclusive. These subsets differ in fold and input composition as well as hardware.

AGS turns this heterogeneity into a testable allocation rule: saved branch computation must exceed the extra probing cost, and the accepted inputs determine the accuracy change.

\subsection{Scope and Extensions}
The common protocol connects system comparison, paired component analysis, and pixel-level failure cases. Query-level aggregation captures dependence among references, while image-cluster uncertainty retains shared query images. Domain transfer tests whether the observed trade-offs extend to specialized imagery, and tiny-reference failures motivate foreground-preserving preprocessing. The central extension question is whether agreement predicts low-value search consistently across domains.

\section{Conclusion}
Visual correspondence and concept guidance yield distinct accuracy--cost trade-offs in training-free in-context segmentation. CG-ICS is more accurate on COCO and LVIS, while INSID3 is substantially faster. Component analysis attributes useful gains to feature processing, semantic concepts, and visual prompts, and reveals a dataset-dependent marginal return from additional search. Reference sensitivity and small-target failures provide further distinctions beyond mean accuracy. These findings motivate AGS, a concrete hypothesis for allocating search through text--visual agreement. Its cost--accuracy formulation specifies the computation savings and accuracy trade-off that an effective adaptive system must achieve.
